\documentclass[letterpaper]{article} % DO NOT CHANGE THIS
\usepackage[submission]{aaai2027}  % DO NOT CHANGE THIS
% The serif, sans-serif, and monospaced fonts are loaded automatically by
% aaai2027.sty. DO NOT add \usepackage{times}, \usepackage{helvet},
% \usepackage{courier}, or any other font package.
\usepackage[hyphens]{url}  % DO NOT CHANGE THIS
\usepackage{graphicx} % DO NOT CHANGE THIS
\urlstyle{rm} % DO NOT CHANGE THIS
\def\UrlFont{\rm}  % DO NOT CHANGE THIS
\usepackage{natbib}  % DO NOT CHANGE THIS AND DO NOT ADD ANY OPTIONS TO IT
\usepackage{caption} % DO NOT CHANGE THIS AND DO NOT ADD ANY OPTIONS TO IT
\frenchspacing  % DO NOT CHANGE THIS
\usepackage{booktabs}
\usepackage{amsmath}
\usepackage{amssymb}   % required for \mathbb (used in Preliminaries)
\pdfinfo{
/TemplateVersion (2027.1)
}

\setcounter{secnumdepth}{2} % section and subsection numbers, used throughout
% for cross-references (e.g. "Section 6.2"); subsubsections not used.

\title{Instruction Tuning Reduces Explanation Overconfidence but Does
Not Consistently Improve Predictive Calibration}
\author{Anonymous Submission}
\affiliations{}

\begin{document}
\maketitle

\begin{abstract} Trustworthy deployment of large language models (LLMs) in high-stakes domains such as medical decision support is often reduced to one fix: calibrate the model's confidence. We ask instead whether a common intervention, instruction tuning, affects \emph{predictive} calibration and \emph{explanation} calibration in the same way. Pairing each base model with its instruction-tuned counterpart across three families (Llama, Qwen, Mistral) $\times$ two QA datasets (MedQA, StrategyQA) $\times$ two explanation regimes (post-hoc, joint), plus ECQA post-hoc, gives 13 base/instruct pairs. Instruction tuning improves explanation calibration (ECE$_{\text{exp}}$, the gap between an explanation's conveyed certainty and its answer's correctness) in 12 of 13 pairs (Wilcoxon $p{=}0.0005$; sign test $p{=}0.003$), yet its effect on predictive calibration is directionally inconsistent (5 of 13; $p{=}0.50$). A difference-of-differences interaction test confirms the contrast (10 of 12, $p{=}0.007$; 95\% CI $[+0.13,+0.35]$), while an equivalence test on the predictive side fails, so we claim inconsistency, not absence. On the two multiple-choice datasets sharing one confidence estimator, instruction tuning \emph{worsens} predictive ECE in 7 of 9 pairs while improving ECE$_{\text{exp}}$ in all 9. All 30 cells are overconfident ($\overline{S} > $ accuracy), and the gain is driven by reduced stated certainty rather than accuracy: $\Delta\overline{S}$, which has no accuracy term, improves in 13 of 13 pairs ($p{=}0.0002$). The finding is therefore reduced explanation \emph{overconfidence}, contingent on an overconfident regime. Decomposing ECE$_{\text{exp}}$ shows this is substantially a shift in \emph{average} certainty toward accuracy (calibration-in-the-large) rather than finer bin-level resolution under the justificatory elicitation. Varying only the prompt shows this degeneracy is a property of the prompt, not the dataset: an assessment framing recovers substantial resolution ($22$--$38\%$ on MedQA, $32$--$70\%$ on ECQA) for all three instruction-tuned models and neither usable base model. An independent model-level analysis corroborates the dissociation: predictive and explanation calibration show no significant correlation across models (Spearman $\rho{=}0.22$, $n{=}11$, $p{=}0.52$). Together these results go beyond a correlational null: a single common training intervention selectively improves one trustworthiness signal while leaving the other model-dependent and inconsistent in direction. Predictive and explanation calibration therefore do not share a single upstream representational cause, and must be measured and intervened on independently. \end{abstract}



\section{Introduction} Trustworthy deployment of LLMs is often treated as separate problems studied by separate literatures. One line of work shows that instruction tuning and RLHF frequently \emph{degrade} predictive calibration even as they improve instruction-following \citep{zhu2023calibration,tian2023justask,leng2025taming}. A second line of work shows that LLM explanations are frequently \emph{unfaithful}: they can be swapped, biased, or manipulated without the stated reasoning changing to reflect it \citep{turpin2023language,atanasova2023faithfulness,lanham2023measuring}. Closest to our own question, a third line of work asks specifically whether models can express their own uncertainty faithfully in words, and finds they largely cannot: verbalized hedging is a poor match for a model's own intrinsic confidence \citep{kadavath2022language,lin2022teaching,yona2024llms}. That literature checks \emph{self-consistency}, verbalized certainty against the model's own confidence, with no reference to whether the underlying answer is actually correct. We ask a calibration-specific question instead: whether explanation certainty is calibrated against \emph{ground-truth correctness}, in the classical \citet{guo2017calibration} sense. This is a distinct and, for a user relying on an explanation, more decision-relevant target, since what matters to them is whether stated certainty predicts correctness, not merely whether it matches the model's own internal state (Section~\ref{sec:prelim-note} makes this distinction formal). None of these three literatures asks the question their combination raises. If instruction tuning changes predictive calibration, does it change explanation calibration too, and in the same direction, as if both are symptoms of one underlying property, or are the two dissociable? If calibration is something a training intervention can target directly, is there any reason to expect it moves explanation quality along with it, or should the two be measured and fixed independently? We test this directly for LLMs. We use instruction tuning specifically because it is a real, common intervention with documented effects on predictive calibration, unlike a mechanical rescaling such as temperature scaling, whose explanations are generated from realized answer text and would show a null effect by construction rather than by finding. Our test spans three model families, a clinical domain (MedQA), and two non-medical commonsense-reasoning domains (StrategyQA, ECQA), so the result is not an artifact of one dataset. If predictive calibration and explanation calibration were downstream of one shared representational quantity, the same intervention acting on that quantity should move both in a consistent, correlated way across model families. We test whether it does.

\paragraph{Contributions.} \begin{enumerate} \item We formalize \textbf{explanation calibration} (ECE$_{\text{exp}}$), classical ECE \citep{guo2017calibration} applied to an explanation's conveyed certainty against ground-truth answer correctness, and distinguish it from self-consistency measures (Section~\ref{sec:prelim-note}). \item Using a direct paired interaction test, we show instruction tuning improves explanation calibration but not predictive calibration across three model families and three datasets; an accuracy-independent measure ($\Delta\overline{S}$, 13/13 pairs) confirms the effect is reduced \emph{overconfidence}, and an independent cross-model analysis finds no correlation between the two signals (Sections~\ref{sec:interaction}, \ref{sec:results}). \item We show the apparent MedQA degeneracy is a property of the elicitation \emph{prompt}, not the dataset: an assessment framing recovers explanation resolution for instruction-tuned models, replicating on ECQA (Section~\ref{sec:framing}). \item We apply uniform exclusion criteria for confidence-estimator degeneracy and include a multiple-choice-only comparison that removes an estimator confound, so the headline result is not built on unreliable cells (Sections~\ref{sec:exclusion}, \ref{sec:results}). \end{enumerate}

\section{Related Work} \paragraph{Calibration and the effect of post-training.} \citet{guo2017calibration} establish that modern neural networks are frequently overconfident and introduce ECE as the standard measure of confidence-accuracy alignment. For LLMs specifically, instruction tuning and RLHF are repeatedly shown to degrade predictive calibration relative to base models even as they improve instruction-following \citep{desai2020calibration,zhu2023calibration,tian2023justask,leng2025taming}, and \citet{geng2024survey} and \citet{xiong2024llms} survey the resulting landscape of calibration and confidence-elicitation methods. We adopt classical ECE unmodified on the predictive side; our contribution is the paired, model-matched design comparing the \emph{same} intervention across two signals, not a new metric or mitigation. \paragraph{Explanation faithfulness.} A separate literature shows LLM explanations frequently fail to reflect the process that produced the answer: \citet{turpin2023language} show stated reasoning does not update to reflect answer-flipping biases the model is demonstrably sensitive to; \citet{atanasova2023faithfulness} and \citet{lanham2023measuring} propose counterfactual and intervention-based tests that find explanations are often unfaithful by these criteria. This work asks whether explanations are faithful to the model's \emph{process}; we ask a narrower, calibration-specific question, whether an explanation's \emph{conveyed certainty} matches ground-truth answer correctness in aggregate, which is neither implied by nor a substitute for faithfulness, but is answerable with the same classical-ECE machinery used on the predictive side. \paragraph{Uncertainty communication.} \citet{kadavath2022language} show LLMs have some capacity to self-evaluate their own claims (e.g., via a learned P(IK) probe), a different question from whether generated text communicates that self-knowledge. \citet{lin2022teaching} fine-tune GPT-3 to emit calibrated word-based confidence, showing faithful verbalized uncertainty is achievable with targeted training, not a free byproduct of instruction tuning alone. \citet{yona2024llms} test whether models do this spontaneously and find they largely do not: verbalized decisiveness is a poor match for a model's own intrinsic confidence across a range of aligned LLMs. All three ask a \emph{self-consistency} question, does expressed certainty match the model's own internal state, not a \emph{calibration} question in the classical \citet{guo2017calibration} sense: whether expressed certainty predicts actual correctness. We ask the latter, for explanation-embedded certainty specifically, at the level of aggregate model behavior rather than per-item self-consistency tracking. \paragraph{Semantic uncertainty estimation.} For StrategyQA's open-ended answer format, where no fixed option set exists to restrict a softmax over, we build on the semantic-entropy line of work \citep{kuhn2023semantic,farquhar2024detecting}, adapting it into an embedding-based clustering estimator over repeated samples (Section~\ref{sec:conf-extraction}). For the multiple-choice datasets (MedQA, ECQA) we instead use a closed-form logit-based estimator that requires no sampling.

\section{Preliminaries} \label{sec:definitions} Let $f_\theta$ be an LLM answering question $x$ with answer $\hat y$, predictive confidence $\text{conf}(x)$, and generated explanation $e$. \paragraph{Predictive ECE.} Standard Expected Calibration Error \citep{guo2017calibration}: items are sorted into $B$ confidence bins, and \[ \text{ECE} = \sum_{b=1}^{B} \frac{n_b}{n} \left| \overline{\text{conf}}_b - \overline{\text{acc}}_b \right|. \] \paragraph{Explanation certainty and explanation ECE.}\label{sec:prelim-note} An LLM judge rates the uncertainty conveyed by $e$'s language, $u(e)\in[0,100]$; we define certainty $S(e) = 1-u(e)/100$ and compute \[ \text{ECE}_{\text{exp}} = \sum_{b=1}^{B} \frac{n_b}{n} \left| \overline{S}_b - \overline{\text{acc}}_b \right|, \] where $\text{acc}$ is, as in predictive ECE above, whether the \emph{answer} $\hat y$ is correct, not any judgment of the explanation's own content. For MedQA and StrategyQA, which constitute the primary grid, no reference explanations exist, so answer correctness is the only verifiable target available. ECQA, used in the paired comparison and the framing experiment (Sections~\ref{sec:pairing}, \ref{sec:framing}), \emph{does} ship human-written reference explanations \citep{aggarwaletal2021ecqa}. We deliberately do not use them in ECE$_{\text{exp}}$ here: scoring explanation content against a reference introduces a second, judge-dependent correctness variable whose construction (a continuous agreement score versus a thresholded correct/incorrect label) materially changes what the resulting quantity measures, and settling that is a separate contribution rather than something to decide midway through this one. Holding ECE$_{\text{exp}}$ fixed to answer correctness also keeps the ECQA framing result directly comparable to the MedQA one. We return to this in Limitations. This is identical in form to predictive ECE with $S(e)$ substituted for $\text{conf}(x)$. We distinguish this explicitly from a hypothetical self-consistency measure that would instead compare $S(e)$ to the model's own $\text{conf}(x)$, with no reference to ground truth at all; call this ExCE (Explanation-Confidence Consistency Error) for contrast. ECE$_{\text{exp}}$ is a calibration measure in the classical sense, checked against ground-truth answer correctness; ExCE would only be a self-consistency measure, checked against nothing external to the model, and the two are not interchangeable: a model could score well on ExCE (explanation certainty matches its own stated confidence) while both are simultaneously miscalibrated against reality, or score poorly on ExCE while both happen to be individually well-calibrated.

\section{Methodology} \label{sec:methodology} This section gives the concrete procedure behind every number in Section~\ref{sec:results}, so that each measure in Section~\ref{sec:definitions} is fully reproducible from raw model outputs.

\subsection{Confidence Extraction} \label{sec:conf-extraction} \paragraph{MedQA.} A single forward pass is run over the question and its five options. At the token position immediately preceding the answer, we read the model's raw logits for exactly the five option-letter tokens $z_A,\ldots,z_E$ and restrict the softmax to this set: \[ q_\ell = \frac{\exp(z_\ell)}{\sum_{\ell'\in L}\exp(z_{\ell'})}, \qquad L=\{A,B,C,D,E\}. \] The predicted answer is $\hat y = \arg\max_\ell q_\ell$ and $\text{conf}(x) = \max_\ell q_\ell$. No sampling is required; this is a deterministic function of one forward pass. 

\paragraph{StrategyQA.} No fixed answer set exists for StrategyQA's open-ended yes/no format, so no comparable softmax restriction is possible. We draw $k{=}10$ temperature-sampled completions ($T{=}0.7$), embed each with \texttt{sentence-transformers/all-MiniLM-L6-v2}, and cluster them via agglomerative clustering with a fixed cosine-distance threshold of $0.15$ and average linkage, giving $m \leq k$ clusters. Let $p_i = n_i/k$ be the fraction of the $k$ samples falling in cluster $i$ (of size $n_i$). We compute Shannon entropy over the resulting cluster-frequency distribution, \[ H(x) = -\sum_{i=1}^{m} p_i \log p_i, \] and normalize by the maximum entropy achievable with $k$ samples, $\log k$, so that $\hat U_{\text{EBSE}}(x) = H(x) / \log k \in [0,1]$ regardless of how many clusters are realized. We define $\text{conf}(x) = 1 - \hat U_{\text{EBSE}}(x)$: more self-agreement across the 10 samples yields higher confidence. \textbf{Note:} we reconstruct this normalization ($\log k$, the theoretical maximum given a fixed sample budget, rather than $\log m$, the realized cluster count) from standard convention in the semantic-entropy literature; confirm it matches the actual implementation before treating this as verified. \paragraph{ECQA.} Identical in form to MedQA: a single forward pass, logits read at the token position immediately preceding the answer, softmax restricted to the option-letter tokens present for that item (ECQA is uniformly 5-option, $L{=}\{A,B,C,D,E\}$), and $\text{conf}(x) = \max_\ell q_\ell$. No new estimator is introduced; this is the same closed-form procedure applied to a second multiple-choice dataset, which is what permits the multiple-choice-only paired comparison in Section~\ref{sec:pairing}. \subsection{Correctness Labels} \paragraph{MedQA and ECQA.} The predicted option letter is extracted from the generated text (regex-based pattern matching for instruction-tuned models; a content-overlap fallback against the option strings for base models, which frequently state the substance of an option rather than its letter) and compared directly to the labeled correct option. No judge is required for either dataset. For ECQA, the correct option is identified by exact string match between the dataset's answer field and one of the five option strings; a row where no exact match is found fails at generation time and is recorded as missing under the same missing-value handling as any other dropped row (Section~\ref{sec:exclusion}), not filtered out in advance. \paragraph{StrategyQA.} Because free-text yes/no answers can be correct under paraphrase, correctness is determined by an independent LLM judge comparing the greedily decoded answer against the dataset's reference label(s), allowing for paraphrase and hedged phrasing rather than requiring exact string or lexical-overlap match. \subsection{Explanation Elicitation and Certainty Scoring} \label{sec:elicitation} \paragraph{Generation regime.} MedQA and StrategyQA use two conditions: \emph{post-hoc}, where the explanation is generated in a separate forward pass conditioned only on the question and the already-realized answer text; and \emph{joint}, where answer and explanation are generated in a single continuous decoding pass. For post-hoc, the base-model prompt is \texttt{"Q: \{question\}\textbackslash nA: \{answer\_letter\}. \{answer\_text\}\textbackslash nThis is because"}; the instruction-tuned prompt is \texttt{"Question: \{question\}\textbackslash n\textbackslash nAnswer: \{answer\_letter\}. \{answer\_text\}\textbackslash n\textbackslash nProvide an explanation for this answer:"}. The joint condition differs structurally, not just in wording: rather than two separate forward passes (one to fix the answer via the closed-form confidence procedure of Section~\ref{sec:conf-extraction}, a second, freshly conditioned on that fixed answer, to elicit the explanation), it issues a single generation call that produces the answer statement and its explanation together in one continuous decoding pass, using the same base/instruct prompt-style distinction as post-hoc. Predictive confidence is computed identically to post-hoc in both conditions, via the separate closed-form logit read described in Section~\ref{sec:conf-extraction}; the joint condition's continuous generation is the source of the explanation text and, where its self-reported answer differs from the confidence-step's greedy answer, of the correctness label as well. \textbf{We do not reproduce the literal joint-condition prompt strings here}: this draft was finalized without direct access to that script, only a structural description of what it does, and we flag this explicitly as a completeness gap rather than an equivalence we have verified character-for-character against post-hoc's templates; the exact templates should be added before this is considered fully reproducible. ECQA uses post-hoc generation only, with the same base/instruct post-hoc templates given above; we did not implement a joint-generation variant for it. \paragraph{Elicitation framing.} Independently of generation regime, we additionally vary the \emph{instruction} under which the post-hoc explanation is elicited, holding the answer and its confidence fixed: \emph{justificatory} (the post-hoc prompts given above, used throughout the primary grid) versus \emph{assessment}, which asks the model to assess rather than justify. For instruction-tuned models the assessment prompt is \texttt{"Question: \{question\}\textbackslash n\textbackslash nAnswer: \{answer\_letter\}. \{answer\_text\}\textbackslash n\textbackslash nAssess whether this answer is correct, and explain your assessment. If you are uncertain whether the answer is correct, say so explicitly and explain why."} For base models we use \texttt{"Q: \{question\}\textbackslash nA: \{answer\_letter\}. \{answer\_text\}\textbackslash nIs this answer correct? I am"} rather than a direct instruction: an early pilot found a bare instruction-style prompt produced empty completions from some base models at a high rate, consistent with a noun-phrase-plus-colon shape being an unnatural continuation point for a model with no chat-turn training; the sentence-fragment framing above resolves this and is used for every base model reported here, including on ECQA. Both framings are run for all six models on ECQA (Section~\ref{sec:framing}); for MedQA, justificatory explanations come from the primary grid, and assessment explanations are drawn from earlier exploratory work by the authors using the same items, models, and assessment-style prompt, re-scored here with the judge described below for direct comparability. Each resulting explanation, under either framing, is scored by an independent LLM judge (GPT-4/5-class) prompted to rate, on a $0$--$100$ scale, how much uncertainty the explanation's \emph{language} conveys to a reader; the judge is not asked to infer the model's latent predictive state, only to read the surface text. We validated this judge against 100 human-annotated explanations, obtaining Spearman $\rho{=}0.709$ ($p<0.001$); that validation was performed on justificatory-framing MedQA/StrategyQA explanations, and we do not re-validate it for assessment-framing or ECQA explanations in the present study (Limitations). \subsection{Calibration Computation} \label{sec:calib-computation} Given $n{=}200$ pairs $(\text{conf}_i, \text{correct}_i)$ (or $(S_i,\text{correct}_i)$ for ECE$_{\text{exp}}$) per model/dataset/condition cell, ECE is computed as follows. \begin{enumerate} \item Sort all $n$ confidence values. \item Partition into $B{=}10$ \emph{adaptive} bins of equal item count (quantile binning), rather than equal-width bins: with confidence distributions that are frequently skewed toward one end of $[0,1]$ (Section~\ref{sec:exclusion}), equal-width bins can leave several bins nearly empty and unable to support a stable estimate. \item Within each bin $b$, compute the mean confidence $\overline{\text{conf}}_b$ and the empirical accuracy $\overline{\text{acc}}_b$ (fraction of items with $\text{correct}{=}1$). \item ECE is the item-count-weighted mean absolute gap between the two across all bins (Section~\ref{sec:definitions}). \end{enumerate} ECE$_{\text{exp}}$ follows the identical procedure with $S_i$ in place of $\text{conf}_i$ (Section~\ref{sec:definitions}).

\subsection{Statistical Inference} \label{sec:stat-inference} Each test is introduced once here; later sections report numbers without re-justifying it. \paragraph{Paired tests.} Every base-vs-instruct comparison (Sections~\ref{sec:pairing}, \ref{sec:interaction}) uses the \textbf{Wilcoxon signed-rank test}, which assumes no normality of the paired differences, appropriate at these sample sizes ($n{=}9$--$16$) where normality cannot be checked. Alongside it we report the more conservative \textbf{sign test}, which uses only the direction of each difference (against a Binomial$(n,0.5)$ null) and so is not driven by one or two large-magnitude pairs; where the two disagree (always flagged), the direction is real by Wilcoxon but not uniform enough for the stricter test, and we report both rather than the more favorable one alone. \paragraph{Equivalence testing (TOST).} A non-significant Wilcoxon result shows only that an effect was not detected, not that it is absent. To test absence directly we use two one-sided Wilcoxon tests against a pre-specified bound $\pm\delta$ ($H_0: \text{median}(d_i)\leq-\delta$ and $H_0: \text{median}(d_i)\geq\delta$); the equivalence $p$-value is the larger of the two, and equivalence within $\pm\delta$ holds only if it falls below $0.05$. We report $\delta\in\{0.05,0.10\}$ ECE points, small relative to the $0.1$--$0.5$ ECE values observed throughout, so passing at either bound would support a genuine absence claim. \paragraph{Correlation.} The model-level relationship between predictive ECE and ECE$_{\text{exp}}$ (Section~\ref{sec:results}) uses \textbf{Spearman} $\rho$ (monotonic, no normality assumption, appropriate for bounded skewed values at $n{=}11$), reported as the cross-sectional analogue of the intervention-level paired test. Because both measures are functions of accuracy by construction, we also report a \textbf{partial correlation} controlling for per-cell accuracy. \paragraph{Effect sizes and robustness.} For each significant paired result we report the \textbf{Hodges--Lehmann} estimator (the median of pairwise averages the Wilcoxon test is built around, robust to outliers, in ECE units). ECE 95\% CIs come from a nonparametric \textbf{bootstrap} (1{,}000 item-level resamples); correlation CIs use the \textbf{Fisher $z$-transform} ($\text{se}=1/\sqrt{n-3}$, approximate at $n{=}11$ and read as indicative). We assess correlation stability with a \textbf{leave-one-model-out jackknife}, recomputing on the remaining five to check the sign and rough magnitude do not depend on a single model.

\subsection{Paired Base/Instruction-Tuned Comparison} \label{sec:pairing} For each of the three model families and each of the four MedQA/StrategyQA dataset$\times$condition combinations, we pair the base model's cell with its instruction-tuned counterpart's cell on the same measure, giving 12 possible pairs. A pair is excluded if either side is flagged under the exclusion criteria of Section~\ref{sec:exclusion}, since a comparison against an unreliable estimate is not meaningful; this leaves 10 of 12 for predictive ECE and 10 of 12 for ECE$_{\text{exp}}$ (different pairs are dropped for each measure, since the two measures have different excluded cells). ECQA contributes one further pair per family (post-hoc only, no excluded cells identified, though we note in Limitations that we did not apply the same severe-bin-collapse check to it that we applied to the primary grid), giving 13 pairs in total for each measure. We test each measure's paired differences (base $-$ instruct) with the Wilcoxon signed-rank test defined above, at these sample sizes ($n{=}10$ for the primary grid, $n{=}13$ with ECQA included).

\subsection{Correcting for Non-Independence Across Conditions} \label{sec:independence} Post-hoc and joint cells for the same model and dataset are not independent observations of that model's calibration: both are computed from closely related generations of the same 200 questions, and their ECE values are frequently near-duplicates. Treating all 20--22 cells as independent in a single pooled Spearman correlation therefore pseudoreplicates each model roughly twofold. Our primary correlation analysis (Table~\ref{tab:correlations}) instead averages predictive ECE and ECE$_{\text{exp}}$ within each model$\times$dataset pair across the two conditions (excluding a condition individually if it is flagged under Section~\ref{sec:exclusion}, and excluding the model$\times$dataset pair entirely from a given correlation if both conditions are flagged), yielding one independent observation per model$\times$dataset combination ($n{\leq}12$). We also report the same three correlations computed separately within each dataset, without averaging across conditions, as a secondary check; these remain pseudoreplicated within dataset but at least do not mix MedQA's closed-form confidence construct with StrategyQA's sampling-based one in a single correlation (Section~\ref{sec:limitations-estimator}). The paired test in Section~\ref{sec:pairing} is unaffected by this issue, since it compares matched cells rather than pooling across models. 

\subsection{Data Quality Validation and Exclusion} \label{sec:exclusion} Three checks are applied uniformly to every cell before any number is reported: \paragraph{Provenance tagging.} Dataset and condition labels are derived from which directory a file was found in, not trusted from a column that may be absent or use inconsistent naming across pipeline stages; where a column does exist, it is validated against the directory-derived label and any disagreement halts processing rather than silently resolving in either direction. This guards against silently pooling post-hoc and joint data, or pooling across datasets, through a labeling mismatch. 

\paragraph{Missing-value handling.} Rows with a missing confidence, correctness, or judge-uncertainty value are dropped from the specific computation they would enter (predictive ECE and ECE$_{\text{exp}}$ are filtered independently, since a row can be valid for one and missing for the other), and the number and proportion dropped is recorded alongside the result rather than silently absorbed. 

\paragraph{Estimator-degeneracy exclusion.} A cell is excluded from a given measure if adaptive binning collapses to $\leq 2$ bins (near-zero resolution, typically caused by a large majority of items sharing near-identical confidence) or if more than $50\%$ of required values are missing. Applying this criterion identifies two predictive-ECE cells and two ECE$_{\text{exp}}$ cells for exclusion (Section~\ref{sec:results}); in the most severe case, 92--95.5\% of items in a cell share a (near-)identical confidence value, indicating near-total answer convergence rather than a usable calibration signal.

\section{Experimental Setup} \paragraph{Data and models.} Three question-answering datasets: MedQA (5-option clinical board exam questions), StrategyQA (binary commonsense reasoning), and ECQA (5-option general-knowledge commonsense reasoning), following \citet{jin2021disease,geva2021did,aggarwaletal2021ecqa}. MedQA and ECQA share a 5-option multiple-choice format and the same closed-form confidence estimator; StrategyQA's open-ended format requires a different, sampling-based one (Section~\ref{sec:conf-extraction}). Six models throughout: Llama-3.1-8B, Qwen2.5-7B, Mistral-7B, each in base and instruction-tuned form \citep{grattafiori2024llama,qwen2025technical,jiang2023mistral}. \paragraph{Scale and coverage across datasets.}Coverage is not uniform across datasets. MedQA and StrategyQA are run in both post-hoc and joint regimes ($200$ items per cell, $6\times2\times2=24$ cells for each measure). This is the only grid the correlation analysis (Table~\ref{tab:correlations}) uses, since that analysis averages each model$\times$dataset pair across its two conditions (Section~\ref{sec:independence}) and ECQA has only one. ECQA is run in the post-hoc regime only, $200$ test items per model, under \emph{both} elicitation framings (justificatory and assessment, Section~\ref{sec:elicitation}), 6 models $\times$ 2 framings $\times$ 200 items. All 24 primary-grid cells and both ECQA framings use the same uncertainty judge, so ECE$_{\text{exp}}$ is computed identically everywhere it is reported. ECQA enters the paired comparison and interaction test (Sections~\ref{sec:pairing}, \ref{sec:interaction}: 3 of 13 pairs) and the elicitation-framing comparison (Section~\ref{sec:framing}), but not the correlation analyses or the offset/resolution decomposition (Table~\ref{tab:decomposition}), which we did not rerun on ECQA under justificatory elicitation alone; Table~\ref{tab:framing} reports ECQA's resolution directly for both framings instead. \paragraph{Reused MedQA assessment-framing data.} The assessment-framing side of Section~\ref{sec:framing}'s MedQA comparison uses explanations generated in earlier exploratory work by the authors, using the same items and models and the assessment-style prompt given in Section~\ref{sec:elicitation}, but scored there only with a different judge assessing rhetorical support rather than conveyed uncertainty; we re-score them here with the uncertainty judge used throughout this paper, so ECE$_{\text{exp}}$ under both framings is computed by identical means. Judge coverage varies across these six files ($104$--$195$ valid rows of $200$; Table~\ref{tab:framing} note). ECQA's assessment-framing data is generated fresh for this paper (Section~\ref{sec:elicitation}), not reused.

\section{Results} \label{sec:results} \subsection{Instruction Tuning Selectively Improves Explanation Calibration} \label{sec:paired-results} Table~\ref{tab:paired} pairs every base model with its instruction-tuned counterpart across all four dataset$\times$condition combinations, for both measures (Section~\ref{sec:pairing}). Within this primary grid, instruction tuning improves ECE$_{\text{exp}}$ in 9 of 10 valid pairs (Wilcoxon signed-rank $p{=}0.004$); the single reversal is Qwen2.5-7B on StrategyQA, joint condition, and it is small (0.209 to 0.221). By contrast, instruction tuning improves predictive ECE in only 5 of 10 valid pairs, an exact split ($p{=}0.85$): it helps every Llama pair and one Mistral pair, but consistently \emph{hurts} Mistral on MedQA and Qwen in every pair tested. (This ECE$_{\text{exp}}$ improvement is substantially a shift in average certainty toward average accuracy rather than finer bin-level resolution, particularly on MedQA; we decompose and quantify this directly in Section~\ref{sec:discussion}.) \begin{table*}[t] \centering \small \begin{tabular}{lllcccc} \toprule Family & Dataset & Condition & \multicolumn{2}{c}{Predictive ECE} & \multicolumn{2}{c}{ECE$_{\text{exp}}$} \\ & & & Base & Instruct & Base & Instruct \\ \midrule Llama & MedQA & post-hoc & 0.185 & 0.151 & 0.467 & 0.283 \\ Llama & MedQA & joint & 0.238 & 0.140 & 0.475 & 0.192 \\ Llama & StrategyQA & post-hoc & 0.443 & 0.117 & 0.348 & 0.179 \\ Llama & StrategyQA & joint & 0.453 & 0.196 & 0.341 & 0.171 \\ \addlinespace Mistral & MedQA & post-hoc & 0.161 & 0.446 & 0.563 & 0.338 \\ Mistral & MedQA & joint & 0.168 & 0.451 & 0.495 & 0.380 \\ Mistral & StrategyQA & post-hoc & 0.503 & 0.259 & --$^{\ddagger}$ & 0.228 \\ Mistral & StrategyQA & joint & 0.490 & --$^{\dagger}$ & --$^{\ddagger}$ & 0.218 \\ \addlinespace Qwen & MedQA & post-hoc & 0.122 & 0.376 & 0.308 & 0.265 \\ Qwen & MedQA & joint & 0.098 & 0.336 & 0.276 & 0.208 \\ Qwen & StrategyQA & post-hoc & 0.149 & 0.250 & 0.236 & 0.194 \\ Qwen & StrategyQA & joint & 0.125 & --$^{\dagger}$ & 0.209 & 0.221 \\ \bottomrule \end{tabular} \caption{Predictive ECE and ECE$_{\text{exp}}$, base vs.\ instruction-tuned, paired by family/dataset/condition. --: pair excluded (Section~\ref{sec:pairing}). $\dagger$/$\ddagger$: same exclusions as Table~\ref{tab:ece_full}. Primary grid only; ECQA pairs are reported in Table~\ref{tab:paired_tests}. \textbf{Predictive ECE}: 5/10 valid pairs improve under instruction tuning, Wilcoxon $p{=}0.85$. \textbf{ECE$_{\text{exp}}$}: 9/10 valid pairs improve, Wilcoxon $p{=}0.004$.} \label{tab:paired} \end{table*} The effect size for ECE$_{\text{exp}}$ is large: matched-pairs rank-biserial $r{=}0.964$, with a Hodges--Lehmann median paired improvement of $0.126$ ECE points (95\% bootstrap CI $[0.044, 0.197]$, excluding zero). Predictive ECE shows the opposite pattern: $r{=}-0.091$ (near zero), median paired difference $-0.013$ (95\% CI $[-0.261, 0.244]$, comfortably spanning zero). Figure~\ref{fig:paired} shows this directly: nearly every ECE$_{\text{exp}}$ line descends from base to instruct, while predictive ECE lines cross in every direction. \begin{figure*}[t] \centering \includegraphics[width=\textwidth]{fig2_paired_slopes.pdf} \caption{Base $\rightarrow$ instruction-tuned comparison, faceted by dataset $\times$ condition so each panel contains exactly one line per model family (color-consistent across panels); no panel mixes multiple lines of the same color. A family missing from a panel is labeled explicitly (excluded per Section~\ref{sec:exclusion}) rather than silently absent. Top row: explanation calibration improves in 9/10 valid panels. Bottom row: predictive calibration shows no consistent direction (5/10 improve).} \label{fig:paired} \end{figure*} This is stronger evidence for dissociation than an absence of correlation: a single intervention applied to every model moves the two signals differently. Were both downstream of one shared representational quantity, it should not reliably improve one while leaving the other's direction unpredictable across families. \subsection{Testing the Dissociation Directly, and Ruling Out Accuracy as the Driver} \label{sec:interaction} Section~\ref{sec:paired-results} establishes the two effects with separate tests. Concluding a dissociation from ``one significant, one not'' would be the difference-between-significant-and-non-significant fallacy, so we test the contrast directly. Both measures are in ECE units, so for each pair valid on \emph{both} we compute $d_i = \Delta\text{ECE}_{\text{exp}} - \Delta\text{ECE}$ (each $\Delta$ the base $-$ instruct paired difference defined in Section~\ref{sec:pairing}) and test whether it differs from zero. We report all paired tests on the full set of 13 base/instruct pairs: 10 from the MedQA/StrategyQA grid plus 3 from ECQA (one per family, post-hoc). Unlike the correlation analysis, the paired tests need only predictive ECE and ECE$_{\text{exp}}$, so ECQA enters here despite having a single condition (Section~\ref{sec:independence}). \begin{table}[h] \centering \small \begin{tabular}{lccc} \toprule Test & $n$ & improve & Wilcoxon $p$ \\ \midrule \multicolumn{4}{l}{\emph{Primary grid only (MedQA, StrategyQA)}} \\ \quad ECE$_{\text{exp}}$ & 10 & 9 & 0.004 \\ \quad Predictive ECE & 10 & 5 & 0.85 \\ \quad $\Delta\overline{S}$ & 10 & 10 & 0.002 \\ \addlinespace \multicolumn{4}{l}{\emph{All three datasets}} \\ \quad ECE$_{\text{exp}}$ & 13 & 12 & \textbf{0.0005} \\ \quad Predictive ECE & 13 & 5 & 0.50 \\ \quad $\Delta\overline{S}$ & 13 & 13 & \textbf{0.0002} \\ \quad Interaction $d_i$ & 12 & 10 & \textbf{0.007} \\ \addlinespace \multicolumn{4}{l}{\emph{Multiple-choice only (MedQA, ECQA)}} \\ \quad ECE$_{\text{exp}}$ & 9 & 9 & 0.004 \\ \quad Predictive ECE & 9 & 2 & 0.027 \\ \quad $\Delta\overline{S}$ & 9 & 9 & 0.004 \\ \bottomrule \end{tabular} \caption{Paired base/instruct tests. ``Improve'' counts pairs where the instruction-tuned model is better. Adding ECQA strengthens every explanation-side result and makes the interaction test significant with a CI excluding zero. The multiple-choice-only rows use a single confidence estimator throughout, removing the MedQA-vs-StrategyQA estimator difference; there, instruction tuning significantly \emph{worsens} predictive ECE (only 2 of 9 pairs improve) while ECE$_{\text{exp}}$ improves in all 9.} \label{tab:paired_tests} \end{table} Table~\ref{tab:paired_tests} reports the results. On all three datasets, ECE$_{\text{exp}}$ improves in 12 of 13 pairs (Wilcoxon $p{=}0.0005$; sign test $p{=}0.003$; Hodges--Lehmann $+0.125$, 95\% CI $[+0.044,+0.183]$), while predictive ECE improves in 5 of 13 ($p{=}0.50$). The interaction test is significant with a confidence interval excluding zero (10 of 12 pairs, Wilcoxon $p{=}0.007$, sign test $p{=}0.039$, Hodges--Lehmann $+0.210$, 95\% CI $[+0.127,+0.352]$): instruction tuning does move the two measures differently, and we no longer rest that claim on comparing two separate tests. Restricting to the two multiple-choice datasets, which share one closed-form confidence estimator and so remove the estimator difference we must otherwise caveat, sharpens the picture: predictive ECE improves in only 2 of 9 pairs, i.e.\ instruction tuning \emph{worsens} it in 7 of 9 (Wilcoxon $p{=}0.027$, though the more conservative sign test is not significant at $p{=}0.18$), while ECE$_{\text{exp}}$ improves in all 9. We also ran the equivalence test flagged as planned in earlier versions of this work: a TOST on the predictive-side paired differences. \textbf{It does not pass}, at either a $0.05$ or a $0.10$ ECE bound ($p{=}0.45$ and $p{=}0.27$ respectively). We can therefore say instruction tuning's effect on predictive calibration is inconsistent in direction, and on multiple-choice data appears negative, but we cannot claim it is absent or negligible; the title and claims of this paper are worded accordingly. \paragraph{Every cell is overconfident, and the effect is a reduction in stated certainty.} The decomposition (Table~\ref{tab:decomposition}) shows ECE$_{\text{exp}}$ is dominated by the mean offset, raising a confound: a model that simply became more accurate would show a smaller offset with no change in explanation behaviour. Reporting the signed offset resolves this.\emph{All 30 cells across the three datasets have $\overline{S} > \text{accuracy}$} (range $+0.076$ to $+0.563$): explanations are overconfident everywhere, in every model, dataset, and condition. Decomposing the paired improvement accordingly, the median reduction in mean stated certainty is $0.108$ while the median accuracy gain is only $0.033$; the certainty shift accounts for roughly $76\%$ of the movement. Most directly, $\Delta\overline{S}$ alone, which does not involve accuracy in any way and is therefore immune to this confound, improves in $13$ of $13$ pairs (Wilcoxon $p{=}0.0002$; sign test $p{=}0.0002$; Hodges--Lehmann median reduction $0.099$, $95\%$ CI $[0.072, 0.134]$, excluding zero). This is the most robust result in the paper, and it is why we state the finding as a reduction in explanation \emph{overconfidence} rather than as an improvement in explanation calibration in the general sense: the improvement is contingent on operating in an overconfident regime, and the same hedging would increase ECE$_{\text{exp}}$ in a well-calibrated or underconfident one. \begin{table}[h] \centering \small \begin{tabular}{llccc} \toprule Dataset & Model (cond.\ averaged) & Acc. & $\overline{S}$ & $\overline{S}-$Acc. \\ \midrule MedQA & Llama-3.1-8B & 0.452 & 0.922 & $+0.471$ \\ MedQA & Llama-3.1-8B-Instruct & 0.568 & 0.805 & $+0.237$ \\ MedQA & Mistral-7B-v0.1 & 0.393 & 0.922 & $+0.529$ \\ MedQA & Mistral-7B-Instruct-v0.3 & 0.448 & 0.807 & $+0.359$ \\ MedQA & Qwen2.5-7B & 0.583 & 0.871 & $+0.287$ \\ MedQA & Qwen2.5-7B-Instruct & 0.575 & 0.811 & $+0.236$ \\ \addlinespace StrategyQA & Llama-3.1-8B & 0.605 & 0.950 & $+0.344$ \\ StrategyQA & Llama-3.1-8B-Instruct & 0.718 & 0.824 & $+0.106$ \\ StrategyQA & Mistral-7B-v0.1 & 0.652 & 0.980 & $+0.328$ \\ StrategyQA & Mistral-7B-Instruct-v0.3 & 0.688 & 0.867 & $+0.180$ \\ StrategyQA & Qwen2.5-7B & 0.760 & 0.964 & $+0.205$ \\ StrategyQA & Qwen2.5-7B-Instruct & 0.653 & 0.842 & $+0.189$ \\ \bottomrule \end{tabular} \caption{Accuracy and mean explanation certainty $\overline{S}$ per model$\times$dataset (averaged over conditions), with the \emph{signed} offset, for the MedQA/StrategyQA grid. All 30 cells across the three datasets (including ECQA, not shown) have $\overline{S}>\text{accuracy}$: explanations are overconfident everywhere. Within every family$\times$dataset pair, the instruction-tuned model has lower $\overline{S}$ than its base counterpart (13/13 valid pairs including ECQA, Wilcoxon $p{=}0.0002$; Section~\ref{sec:interaction}).} \label{tab:signed_offset} \end{table} \subsection{Predictive and Explanation Calibration, Full Grid} Table~\ref{tab:ece_full} reports the full grid underlying Table~\ref{tab:paired}. Predictive calibration varies substantially across models: Qwen2.5-7B (base) and Llama-3.1-8B-Instruct are best-calibrated (ECE $\approx 0.10$--$0.15$ in most cells), while the worst individual cells among Mistral-7B-Instruct and Qwen2.5-7B-Instruct reach ECE $0.451$ and $0.376$ respectively, $4.6\times$ the single best cell in the table ($0.098$, Qwen2.5-7B base, MedQA joint). Post-hoc and joint predictive ECE are broadly similar for most models (within $\sim 0.03$--$0.05$), with the joint condition's excluded cells being the notable exceptions. \begin{table*}[t] \centering \small \begin{tabular}{llcccc} \toprule Dataset & Model & Condition & Predictive ECE & 95\% CI & ECE$_{\text{exp}}$ \\ \midrule MedQA & Llama-3.1-8B & post-hoc & 0.185 & [0.141, 0.263] & 0.467 \\ MedQA & Llama-3.1-8B & joint & 0.238 & [0.185, 0.305] & 0.475 \\ StrategyQA & Llama-3.1-8B & post-hoc & 0.443 & [0.380, 0.521] & 0.348 \\ StrategyQA & Llama-3.1-8B & joint & 0.453 & [0.390, 0.528] & 0.341 \\ \addlinespace MedQA & Llama-3.1-8B-Instruct & post-hoc & 0.151 & [0.121, 0.230] & 0.283 \\ MedQA & Llama-3.1-8B-Instruct & joint & 0.140 & [0.108, 0.221] & 0.192 \\ StrategyQA & Llama-3.1-8B-Instruct & post-hoc & 0.117 & [0.081, 0.185] & 0.179 \\ StrategyQA & Llama-3.1-8B-Instruct & joint & 0.196 & [0.147, 0.260] & 0.171 \\ \addlinespace MedQA & Mistral-7B-v0.1 & post-hoc & 0.161 & [0.126, 0.237] & 0.563 \\ MedQA & Mistral-7B-v0.1 & joint & 0.168 & [0.127, 0.247] & 0.495 \\ StrategyQA & Mistral-7B-v0.1 & post-hoc & 0.503 & [0.434, 0.569] & 0.367$^{\ddagger}$ \\ StrategyQA & Mistral-7B-v0.1 & joint & 0.490 & [0.421, 0.563] & 0.289$^{\ddagger}$ \\ \addlinespace MedQA & Mistral-7B-Instruct-v0.3 & post-hoc & 0.446 & [0.382, 0.514] & 0.338 \\ MedQA & Mistral-7B-Instruct-v0.3 & joint & 0.451 & [0.385, 0.520] & 0.380 \\ StrategyQA & Mistral-7B-Instruct-v0.3 & post-hoc & 0.259 & [0.203, 0.339] & 0.228 \\ StrategyQA & Mistral-7B-Instruct-v0.3 & joint & 0.297$^{\dagger}$ & [0.241, 0.364] & 0.218 \\ \addlinespace MedQA & Qwen2.5-7B & post-hoc & 0.122 & [0.085, 0.190] & 0.308 \\ MedQA & Qwen2.5-7B & joint & 0.098 & [0.078, 0.181] & 0.276 \\ StrategyQA & Qwen2.5-7B & post-hoc & 0.149 & [0.109, 0.212] & 0.236 \\ StrategyQA & Qwen2.5-7B & joint & 0.125 & [0.099, 0.197] & 0.209 \\ \addlinespace MedQA & Qwen2.5-7B-Instruct & post-hoc & 0.376 & [0.309, 0.448] & 0.265 \\ MedQA & Qwen2.5-7B-Instruct & joint & 0.336 & [0.273, 0.404] & 0.208 \\ StrategyQA & Qwen2.5-7B-Instruct & post-hoc & 0.250 & [0.194, 0.316] & 0.194 \\ StrategyQA & Qwen2.5-7B-Instruct & joint & 0.338$^{\dagger}$ & [0.279, 0.403] & 0.221 \\ \bottomrule \end{tabular} \caption{Predictive ECE (with 95\% bootstrap CI) and ECE$_{\text{exp}}$, all 6 models $\times$ 2 datasets $\times$ 2 conditions ($n{=}200$ per cell). $\dagger$: excluded, predictive-ECE quantile binning collapsed to $\leq 2$ bins (near-total answer convergence). $\ddagger$: excluded, ECE$_{\text{exp}}$ built on $<50\%$ valid judge scores with heavy floor mass among the remainder.} \label{tab:ece_full} \end{table*} The full grid also illustrates that predictive calibration is itself strongly dataset-dependent, independent of the instruction-tuning question. Mistral-7B (base) has among the best predictive calibration in the study on MedQA (ECE $0.161$--$0.168$) but the single worst predictive ECE in the entire table on StrategyQA ($0.490$--$0.503$); its explanation calibration is the worst in the table on MedQA (ECE$_{\text{exp}}$ $0.495$--$0.563$) regardless of dataset. This dataset-sensitivity is a property of all three base models in the same direction (StrategyQA predictive ECE exceeds MedQA predictive ECE for every base model) but at very different magnitudes: Llama-base $2.1\times$ worse, Mistral-base $3.0\times$ worse, Qwen-base only $1.25\times$ worse. Part of this gap plausibly reflects the differing confidence estimators rather than the model alone (Section~\ref{sec:limitations-estimator}). \subsection{The Offset-Degeneracy Is Prompt-Dependent, Not a Property of the Dataset} \label{sec:framing} The decomposition in Section~\ref{sec:discussion} (Table~\ref{tab:decomposition}) shows ECE$_{\text{exp}}$ on MedQA is $97$--$100\%$ calibration-in-the-large, essentially no bin-to-bin resolution. Every cell contributing to that finding was generated under a single \emph{justificatory} elicitation prompt (``Provide an explanation for this answer''). In earlier exploratory work we found that elicitation framing substantially changes how much support an explanation expresses; we therefore re-scored the assessment-framing MedQA explanations from that work (``Assess whether this answer is correct\ldots'') with the same uncertainty judge, holding dataset, items, models, and judge fixed and varying only the prompt. To test whether any effect is specific to the clinical domain, we additionally ran the full pipeline, both elicitation prompts and all six models, on ECQA, a non-medical commonsense-reasoning benchmark with the same 5-option structure. \begin{table}[h] \centering \small \begin{tabular}{lcccc} \toprule Model & \multicolumn{2}{c}{MedQA} & \multicolumn{2}{c}{ECQA} \\ & Just. & Assess. & Just. & Assess. \\ \midrule Llama-3.1-8B-Instruct & 0.3\% & 21.6\% & 0.0\% & 32.3\% \\ Mistral-7B-Instruct-v0.3 & 0.0\% & 22.1\% & 19.2\% & 61.1\% \\ Qwen2.5-7B-Instruct & 0.0\% & 37.9\% & 10.6\% & 70.4\% \\ \addlinespace Llama-3.1-8B (base) & 0.0\% & 0.0\% & 0.0\% & 0.2\% \\ Mistral-7B-v0.1 (base) & 0.0\% & 0.0\% & 0.0\% & 0.0\% \\ Qwen2.5-7B (base) & 3.0\% & 20.1\%$^{*}$ & 21.1\% & 11.5\% \\ \bottomrule \end{tabular} \caption{Percentage of ECE$_{\text{exp}}$ attributable to genuine bin-to-bin resolution, under justificatory vs.\ assessment elicitation, on MedQA and on ECQA (both post-hoc). All three instruction-tuned models recover substantial resolution on both datasets; Llama and Mistral base models recover on neither. $^{*}$Unverified: this MedQA cell's judge output dropped 94 of 198 rows for reasons we have not diagnosed, versus $0$--$5$ for every other MedQA cell; we do not treat it as a base-model data point. On ECQA the same model instead shows a modest decrease, so Qwen2.5-7B (base) is anomalous on both datasets without being consistently so.} \label{tab:framing} \end{table} Table~\ref{tab:framing} reports the result on MedQA and, as an independent replication, on ECQA \citep{aggarwaletal2021ecqa}, a non-medical commonsense-reasoning dataset with the same 5-option multiple-choice structure, generated and judged by the identical pipeline. For all three instruction-tuned models on both datasets, resolution rises substantially (MedQA $21.6$--$37.9\%$, ECQA $32.3$--$70.4\%$): under assessment framing, explanation certainty carries real item-level signal about correctness where justificatory framing yielded little or none. Llama-3.1-8B and Mistral-7B-v0.1 (base) remain fully degenerate under both prompts on both datasets. Qwen2.5-7B (base) is the one model that does not fit cleanly: its MedQA cell is unusable (see caption) and its ECQA resolution \emph{falls} slightly under assessment framing, so we do not count it as evidence in either direction. Two consequences follow. First, the offset-degeneracy documented above is \emph{not} a fixed property of MedQA, as a reader of the decomposition alone might reasonably conclude; it is a property of the justificatory elicitation under which those explanations were generated, and it lifts on a second, non-medical dataset under the same intervention. Second, the recovery is specific to instruction-tuned models, which is consistent with, and independent evidence for, this paper's central claim that instruction tuning is the intervention that makes explanation-side signals informative (Section~\ref{sec:paired-results}). Base models given the same assessment prompt do not produce resolved explanation certainty on either dataset, suggesting the capacity to modulate expressed certainty in response to an assessment instruction is itself something instruction tuning confers. \subsection{Predictive vs.\ Explanation Calibration: No Correlation Across Models} Table~\ref{tab:correlations} reports the Spearman correlation between predictive ECE and ECE$_{\text{exp}}$ after averaging each within model$\times$dataset pair across conditions (Section~\ref{sec:independence}), our primary analysis. The correlation is far from significant, and is not rescued by correcting for pseudoreplication; if anything the correction weakens the pooled estimate. \begin{table}[h] \centering \small \begin{tabular}{lcccc} \toprule Comparison & $n$ & $\rho$ & 95\% CI & $p$ \\ \midrule Pred.\ ECE vs.\ ECE$_{\text{exp}}$ & 11 & 0.218 & [$-$0.44, 0.72] & 0.519 \\ \bottomrule \end{tabular} \caption{Spearman correlation, one observation per model$\times$dataset (conditions averaged; Section~\ref{sec:independence}). 95\% CI via Fisher $z$-transform (approximate at this $n$). Not significant; at this sample size we can rule out a strong correlation but not a moderate one (Limitations).} \label{tab:correlations} \end{table} \paragraph{Robustness: exclusions restored.} Recomputing the correlation with the four estimator-degeneracy exclusions (Section~\ref{sec:exclusion}) reversed, treating every one of the 24 cells as valid, leaves it non-significant ($\rho{=}0.273$, $p{=}0.391$). The null in Table~\ref{tab:correlations} is not an artifact of our exclusion criteria. \paragraph{Robustness: leave-one-model-out jackknife.} Dropping each of the six models in turn and recomputing the correlation on the remaining five keeps it within a comparatively narrow band ($\rho \in [0.08, 0.42]$) and it never changes sign, unlike several other exploratory correlations we checked during this analysis and do not report as findings. We read this stability as one further reason to trust the null result, not as evidence against it. \paragraph{Robustness: partial correlation controlling for accuracy.} Both measures are functions of accuracy by construction, so a natural concern is that accuracy itself is a confound generating or suppressing the correlation in Table~\ref{tab:correlations}. Computing it as a partial correlation controlling for per-cell accuracy (averaged across conditions) leaves it essentially unchanged and still null ($\rho{=}0.075$, $p{=}0.837$, down from $0.218$ uncontrolled). Accuracy is not manufacturing or suppressing this particular null. \paragraph{Per-dataset breakdown.} Computed separately within each dataset (without averaging across conditions, so still pseudoreplicated within dataset; Section~\ref{sec:independence}), StrategyQA alone gives $\rho{=}0.643$ ($n{=}8$, $p{=}0.086$); MedQA alone is not even nominally significant. Neither survives at the primary, corrected sample size in Table~\ref{tab:correlations}. We report this for transparency, not as a finding: at $n{\leq}12$ within a single dataset, a single influential model can swing $\rho$ substantially, and we do not think it is trustworthy on its own.

\section{Discussion} \label{sec:discussion} Instruction tuning improves explanation calibration while leaving predictive calibration's direction unpredictable across families (Section~\ref{sec:paired-results}), and the two respond differently to the \emph{same applied intervention} rather than merely failing to co-vary as models differ. For a deploying organization, the shorthand ``we calibrated the model'' therefore licenses no inference about explanation trustworthiness, and an intervention chosen to improve one signal (instruction tuning, RLHF) cannot be assumed to improve, or even leave unharmed, the other. \paragraph{Offset vs.\ resolution: a construct-validity concern we can now confirm directly, not just flag.} ECE$_{\text{exp}}$ checks whether the certainty score $S(e)$ tracks whether the \emph{answer} is correct, not whether the explanation's content is correct. Since $S(e)$ is not a probability the model committed to, nothing forces it onto the same scale as accuracy. We decompose each cell's ECE into \emph{calibration-in-the-large} ($\text{CITL} = |\overline{S}-\overline{\text{acc}}|$, the pure mean-offset component, equal to the ECE obtained from a constant-$S$ baseline with zero bin-to-bin resolution by construction) and a \emph{resolution residual} ($\text{ECE}_{\text{exp}} - \text{CITL}$, the remainder, reflecting genuine bin-to-bin structure). Table~\ref{tab:decomposition} reports the result. On MedQA, ECE$_{\text{exp}}$ is $97$--$100\%$ offset in all 12 cells (mean $0.3\%$ resolution): the measure is functionally equivalent to comparing two population means, with no bin-level structure at all. On StrategyQA the picture is genuinely mixed (mean $12.6\%$ resolution, up to $49\%$ for Llama-3.1-8B-Instruct). This is not an artifact of the decomposition itself: predictive ECE retains substantial resolution on \emph{both} datasets (mean $16.3\%$ MedQA, $13.6\%$ StrategyQA), so the near-total collapse is specific to ECE$_{\text{exp}}$ on MedQA, not a general property of our estimator. \begin{table}[h] \centering \small \begin{tabular}{llcccc} \toprule Measure & Dataset & Mean & Median & Min & Max \\ \midrule Predictive ECE & MedQA & 16.3\% & 0.5\% & 0.0\% & 98.7\% \\ Predictive ECE & StrategyQA & 13.6\% & 0.8\% & 0.0\% & 56.2\% \\ ECE$_{\text{exp}}$ & MedQA & 0.3\% & 0.0\% & 0.0\% & 3.0\% \\ ECE$_{\text{exp}}$ & StrategyQA & 12.6\% & 8.0\% & 0.0\% & 49.1\% \\ \bottomrule \end{tabular} \caption{Percentage of each cell's ECE attributable to genuine bin-to-bin resolution (the remainder is calibration-in-the-large, pure offset), by measure and dataset, $n{=}12$ cells each. All cells here use justificatory elicitation. ECE$_{\text{exp}}$ on MedQA is degenerate under this prompt; Section~\ref{sec:framing} shows this is a property of the elicitation, not of the dataset.} \label{tab:decomposition} \end{table} \begin{figure*}[t] \centering \includegraphics[width=\textwidth]{fig4_decomposition.pdf} \caption{ECE$_{\text{exp}}$ per model, decomposed into offset (blue) and resolution (red), averaged across conditions. MedQA bars are almost entirely offset for all six models; StrategyQA bars show a visible resolution component for four of five valid models (Llama-base is the exception). Mistral-7B (base) is excluded on StrategyQA (Section~\ref{sec:exclusion}).} \label{fig:decomposition} \end{figure*} \paragraph{What this means for our headline result.} The paired Wilcoxon test (Section~\ref{sec:paired-results}) is a statistical fact about ECE$_{\text{exp}}$ values and is unaffected by this decomposition. What it means substantively is narrower than ``explanation calibration improves'' implies at face value: under justificatory elicitation on MedQA, the finding is more precisely that instruction tuning brings a model's \emph{average} explanation certainty closer to its \emph{average} accuracy (calibration-in-the-large), not that its explanations become better resolved across difficulty levels. On StrategyQA, where resolution is often non-trivial, the richer reading is more defensible for at least some models (notably Llama-3.1-8B-Instruct). Critically, this limitation is tied to the elicitation prompt, not to the dataset. Section~\ref{sec:framing} shows that re-eliciting explanations on the same MedQA items under an assessment framing recovers $21.6$--$37.9\%$ resolution for all three instruction-tuned models. The correct statement is therefore not that ECE$_{\text{exp}}$ is degenerate on MedQA, but that it is degenerate \emph{under justificatory elicitation}, and that instruction-tuned models recover resolved explanation certainty when asked to assess rather than to justify. We report the decomposition and this framing comparison together for exactly this reason: the decomposition alone would license a conclusion about the dataset that the framing experiment shows to be wrong. 

\section{Limitations} \label{sec:limitations} 

\paragraph{Scope and statistical power.} The 24-cell primary grid is six models $\times$ two datasets (MedQA, StrategyQA) $\times$ two generation regimes; this is the grid the correlation analysis uses, since it averages each model$\times$dataset pair across its two conditions (Section~\ref{sec:independence}) and ECQA has only one. ECQA cannot enter the correlation analysis for this reason, but it \emph{does} enter the paired comparison and interaction test (Sections~\ref{sec:pairing}, \ref{sec:interaction}), which need only predictive ECE and ECE$_{\text{exp}}$: 3 of the 13 paired base/instruct comparisons come from ECQA. This gives $n{=}11$ for the corrected correlation test (Section~\ref{sec:independence}), $n{=}13$ per paired comparison including ECQA ($n{=}10$ for the primary grid alone; Section~\ref{sec:pairing}), and $n{=}12$ for the interaction test (Section~\ref{sec:interaction}). The 95\% confidence interval on our correlation estimate (Table~\ref{tab:correlations}) spans roughly $-0.4$ to $0.7$: at this sample size we can rule out a \emph{strong} correlation between predictive and explanation calibration, not a moderate one, and we do not claim otherwise. \paragraph{Multiple comparisons.} We designate the paired test on $\Delta\overline{S}$ (Section~\ref{sec:interaction}) and the paired test on ECE$_{\text{exp}}$ (Section~\ref{sec:paired-results}) as primary; both survive Bonferroni correction for two tests ($p{=}0.002$, $p{=}0.004$). The interaction test, TOST, the correlation, the per-dataset breakdown, the jackknife, and the partial correlation are secondary or exploratory, are reported without correction, and should be read accordingly. \paragraph{Reference explanations left unused.} ECQA provides human-written reference explanations for both correct and incorrect options \citep{aggarwaletal2021ecqa}, which would in principle permit a stronger measure than ECE$_{\text{exp}}$: whether an explanation's \emph{reasoning} is correct, not merely whether the answer it accompanies is. We do not use them here. Doing so requires committing to how a judge's continuous agreement score becomes a correctness target (used continuously, or thresholded, and if thresholded, where) and that choice changes what the resulting quantity measures rather than merely how precisely it is estimated. Introducing it for one of three datasets would also make the framing comparison (Section~\ref{sec:framing}) non-comparable across datasets, which is the comparison that carries the result. We regard reference-grounded explanation correctness as a distinct measurement contribution and leave it to separate work. \paragraph{Tests we have not run.} Two checks bearing directly on the construct are absent. An AUROC of $S(e)$ predicting item-level correctness, paired base vs.\ instruct, would test whether explanations become \emph{discriminative} about correctness independently of the mean-offset arithmetic; and a shuffled-label permutation null (recomputing ECE$_{\text{exp}}$ with correctness permuted within cell) would test directly whether $S(e)$ carries item-level signal at all. Both require per-item joins we did not have in place in time. Their absence is why we rest the headline on $\Delta\overline{S}$, which is accuracy-independent by construction, rather than on ECE$_{\text{exp}}$ magnitudes alone. \paragraph{Estimator differences across datasets.} \label{sec:limitations-estimator} MedQA confidence is a closed-form softmax probability over five options; StrategyQA confidence is derived from entropy over ten repeated samples of an open-ended answer space (Section~\ref{sec:conf-extraction}). These are different constructs on conceptually similar but not identical scales, and every base model's StrategyQA predictive ECE is worse than its MedQA predictive ECE, by a magnitude that varies substantially across families ($1.25\times$ to $3.0\times$; Section~\ref{sec:results}). We cannot fully rule out that part of this gap is an estimator property rather than a model property, which is also why our primary correlation analysis (Table~\ref{tab:correlations}) averages within model$\times$dataset rather than pooling datasets together, and why we additionally report, and caution against over-trusting, the per-dataset breakdown. \paragraph{ECE$_{\text{exp}}$ construct validity.} The offset-vs-resolution decomposition (Discussion, Table~\ref{tab:decomposition}) confirms ECE$_{\text{exp}}$ is 97--100\% offset on MedQA for every one of the 12 justificatory cells: under that prompt it is a comparison of two population means, not a resolved calibration curve. The framing comparison (Section~\ref{sec:framing}) shows this is not a fixed property of the measure or the dataset, resolution returns under assessment framing for instruction-tuned models, but it does mean ECE$_{\text{exp}}$ values should not be interpreted as resolved calibration without checking the decomposition for the elicitation used. One caveat on that comparison: one cell (Qwen2.5-7B base) had 94 of 198 rows dropped by the judge for reasons we have not diagnosed, so we do not count it as evidence either way. \paragraph{Natural intervention, not a designed experiment.} Instruction tuning is a real, common intervention, and comparing base to instruction-tuned models lets us ask whether a shared cause moves both trustworthiness signals together, a stronger design than pure cross-sectional correlation. It is not a controlled experiment: we do not control what instruction tuning changes beyond calibration, and a skeptic could argue some other correlated change (e.g., in verbosity or hedging vocabulary specifically) drives the ECE$_{\text{exp}}$ effect rather than calibration-relevant representational change per se. We do not have a way to rule this out with the present data. \paragraph{Exclusions.} Two of 24 predictive-ECE cells and two of 24 ECE$_{\text{exp}}$ cells in the primary grid were excluded for estimator degeneracy (Section~\ref{sec:exclusion}), reducing the primary-grid paired comparison to $n{=}10$ for each measure ($n{=}13$ with ECQA's 3 pairs included; we did not apply the same severe-bin-collapse check to the ECQA cells and flag this as a gap in our own exclusion procedure rather than evidence they are clean). These exclusions are concentrated in specific models and conditions (base models under StrategyQA joint generation), which may itself be a signal worth independent study rather than pure noise. 

\paragraph{Judge dependence.} ECE$_{\text{exp}}$ depends on a single LLM judge's certainty rating, validated against human annotation at $\rho{=}0.709$ (Section~\ref{sec:elicitation}) but not independently re-validated for the present study. \section{Conclusion} Instruction tuning, a single common training intervention, improves explanation calibration in 12 of 13 valid paired comparisons across three model families and three datasets, while its effect on predictive calibration is inconsistent in direction (5 of 13), and is significantly negative on the two multiple-choice datasets considered alone. All 30 cells in the study are overconfident on the explanation side, and the effect is driven mainly by a reduction in stated certainty rather than by accuracy gains: $\Delta\overline{S}$, which does not involve accuracy at all, improves in 13 of 13 pairs ($p{=}0.0002$). The finding is therefore best stated as a reduction in explanation \emph{overconfidence}, contingent on an overconfident regime, rather than as improved explanation calibration in general. A direct interaction test on the difference-of-differences establishes the contrast between the two measures ($p{=}0.007$, CI excluding zero); an equivalence test on the predictive side does not pass, so we claim its effect is inconsistent, not absent. Under the justificatory elicitation used across that grid, this improvement is substantially a shift in average certainty toward average accuracy (calibration-in-the-large) rather than finer bin-level resolution on MedQA (Table~\ref{tab:decomposition}); varying only the elicitation prompt shows that this degeneracy belongs to the prompt rather than the dataset, and that an assessment framing recovers substantial resolution for instruction-tuned models but not for base ones, on MedQA and, independently, on the non-medical ECQA benchmark (Table~\ref{tab:framing}). We report these analyses rather than let the headline pairing result imply a richer form of calibration than the data supports. Corroborating model-level correlational evidence, corrected for non-independence across generation conditions, finds no significant relationship among predictive calibration, explanation calibration, and truth-sensitive explanatory support. Together these indicate that predictive calibration and explanation calibration do not share a single upstream cause: a training intervention chosen to help one property cannot be assumed to help, or even leave unharmed, the other. Trustworthy deployment in domains such as medical decision support requires auditing and, where necessary, separately intervening on each. \section*{Ethical and Broader Impact Statement} This work measures, but does not certify, trustworthiness properties of LLMs, in one clinical (MedQA) and two non-medical commonsense-reasoning (StrategyQA, ECQA) question-answering settings; none of the models studied are deployed in clinical care, and no claims here should be read as endorsing any model's readiness for clinical use. Our central finding is cautionary: it argues against treating calibration as a sufficient safeguard and for independent auditing of explanation behavior before deployment in high-stakes settings.



\bibliography{references}
\end{document}

% ============================================================================
% PAGE BUDGET -- AAAI 2026 limit is 7 pages content + 1 page references only
%
% Current draft: ~6,520 prose words, 5 tables, 3 figures.
% Estimated in AAAI two-column: ~8.3 pages (5.9 prose + 2.4 floats).
% Roughly 1.3 pages / ~1,400 words over, if all 8 floats are kept.
%
% NOTE: this is an estimate from word/float counts, not a real compile --
% get the true number by compiling with aaai2026.sty before deciding cuts.
%
% Cut candidates, roughly in order of least damage to the argument:
%
%  1. Table 2 (full 24-cell grid, tab:ece_full) -> move to appendix.
%     Saves ~0.5 page. Tables 1 and 4 carry the actual arguments; Table 2
%     is reference data. Biggest single win.
%
%  2. Figure 2 (scatterplots, fig:scatter) -> appendix. Saves ~0.3 page.
%     The correlations are null and already fully stated in Table 3; the
%     figure shows absence of a pattern, which the table already conveys.
%
%  3. Per-dataset breakdown paragraph (Section 6) -- explicitly flagged in
%     the text as "not findings"; ~120 words.
%
%  4. Related Work trimming -- typically compressible by 30-40% without
%     losing coverage; ~200-300 words.
%
%  5. The two robustness paragraphs (exclusions-restored, jackknife) could
%     compress to 3-4 sentences total, keeping the numbers and dropping the
%     methodological narration; ~200 words.
%
% Keep regardless: Table 1 (paired result), Table 4 (decomposition),
% Table 5 (framing comparison), Figure 1 (paired slopes), and the
% Limitations section (AAAI reviewers weight it, and its content is
% load-bearing for the honesty of the claims).
% ============================================================================
